\section{Numerical Evaluation}\label{sec:numerics}
The numerical evaluation has two parts. Sections~\ref{sec:constants}--\ref{sec:num-loadaxis}
verify the theory term by term under the known true load and sign composition, while
Section~\ref{sec:e2e} constructs an upper bound from only the quantities observable after
failure and checks its coverage. Table~\ref{tab:exp-overview} lists the question addressed in
each section and the corresponding theoretical result.

\begin{table}[!htbp]
\centering
\caption{Overview of the experiments.}
\label{tab:exp-overview}
\small
\begin{tabular}{lp{.5\linewidth}ll}
\toprule
Section & Question tested & Corresponding theory & Result\\
\midrule
\ref{sec:constants} & the conditional variance, the direction of the conditional mean, and the decay of $\rho^2$ & Theorem~\ref{thm:main} & Figure~\ref{fig:selection}\\
\ref{sec:threshold} & the shift of the threshold center at finite size & the shift constant $C_k$ & Figure~\ref{fig:threshold}\\
\ref{sec:chisq} & the shape of the unconditional measurement under balanced signs & the chi-square correction & Figure~\ref{fig:chisq}\\
\ref{sec:calib} & the accuracy of the lower-tail probability after combining the two corrections & Eq.~\eqref{eq:corrected} & Figure~\ref{fig:calib}\\
\ref{sec:num-loadaxis} & the separation of the macro layer and the repeated-edge layer at low load & Proposition~\ref{prop:lowload}, \ref{prop:mixture} & Figure~\ref{fig:twomech}\\
\ref{sec:e2e} & the coverage and width of the upper bound $U$ when $d$ is unknown & the envelope rule and the mixture rule & Table~\ref{tab:e2e}\\
\bottomrule
\end{tabular}
\end{table}

\subsection{Experimental Design and Comparison Targets}\label{sec:exp-design}
\emph{Simulation setup.} The experiments use an independent Plain mapping and ideal peeling.
The analytic constants are computed by a separate script and are not fitted to the simulation
data. The left-tail series contains $k\in\{3,5\}$, $M\in\{256,1024,4096,16384,65536\}$, and
four target coordinates $\{-2,-2.5,-3,-3.5\}$, for 40 configurations in total. Because $d$
must be an integer, the actual window location is determined by the chosen integer $d$.

\emph{Three coordinates.} The first is the theorem coordinate:
\begin{equation*}
z_0=\frac{(d/M-\alpha_c)\sqrt M}{w_k},\qquad w_k=\sigma_k/|\beta_k|.
\end{equation*}
The second is the shifted coordinate $z_{\rm sh}=z_0+C_kM^{-1/6}$, with the numerical values
of the shift constant $C_k$ in Section~\ref{sec:rules} and the closed form in
Appendix~\ref{app:ck}. The result of DM concerns the configuration model;
Section~\ref{sec:threshold} verifies this shift on the Plain model. A sequence obtained by
fixing $z_{\rm sh}$ corresponds to a single limiting $r$, and therefore satisfies the
conditions of the theorem. The third is the empirical coordinate $z_{\rm emp}=\Phi^{-1}(\widehat p_F)$,
recovered from the measured failure rate and used only for diagnosis.

\emph{Classification of failures.} A run fails when the 2-core is nonempty after peeling.
According to the number of edges in the remaining core, failures are divided into macro-core
failures and small-core failures; a small core containing a pair of identical edges is
assigned to the repeated-edge layer. A macro core is identified by the core edge count
exceeding $0.1d$\footnote{Direct sampling uses the threshold $0.1M$. On all configurations
with $M\ge1024$ the two criteria give the same conclusion; the difference appears only in the
41 configurations with $M\le256$.}.

\emph{Data split.} For configurations with complete records, the samples are split by trial
index into two disjoint halves: the empirical threshold and the empirical failure rate come
from the training half, and the actual lower-tail probability is computed on the failing
samples of the test half. The calibration table contains 624 configuration/sign rows,
involving 115 distinct $(k,M,d)$.

\emph{Thresholds compared.} Sections~\ref{sec:constants}--\ref{sec:num-loadaxis} use the true
load and sign composition and compare four thresholds: the unconditional normal threshold,
the unconditional chi-square threshold, the threshold given by the selection law, and the
corrected selection threshold of Eq.~\eqref{eq:corrected} (with $n=d$ and
$b^2/n^2=\rho^2$). Both corrections come from analytic formulas and contain no fitted
coefficients; the basis for the chi-square term is given in Section~\ref{sec:chisq}.
Section~\ref{sec:e2e} compares the envelope rule and the mixture rule of
Section~\ref{sec:rules}, with the unconditional threshold as the baseline and $d$ given by the
inversion of Section~\ref{sec:upper}.
